\documentclass{article}

\usepackage[T1]{fontenc}
\usepackage{amsmath, amssymb}
\usepackage{ngerman}
\usepackage{graphicx}
\usepackage{tabularx}
\usepackage{mathptmx}
\usepackage[scaled=.92]{helvet}
\usepackage{courier}

\begin{document}

\parindent0cm % verhindert das Einrücken von Absätzen 

% Hier wird die Titelseite gebastelt
% Zuerst wird die Seite auf empty eingestellt, das heißt sie hat keine 
% Seitenzahl und auch keine Kopf- bzw. Fußzeile
\thispagestyle{empty}

% Der Titel wird zentriert in fetter riesiger Schrift gesetzt.

\begin{center}
\begin{Huge}
\textbf{Bachelorprojekt Informatik \\ \vspace{0.2cm}
CS3701-KP05\\ \vspace{2.5cm}
Institut für Telematik\\ \vspace{0.2cm}
Wintersemester 2025/26\\\vspace{2.5cm}
Entwicklung einer Android-App zur mobilen
Überwachung der H0-Modellbahn}\vspace{1.5cm}
\end{Huge}

% Weiter in großer Schrift

\begin{huge}
Thorben Heinz, 756972 \\
<Name>, <Matriklenummer> \\
Lübeck, \today 

\textbf{Leitung:} Prof. Dr. Stefan Fischer

\textbf{Mitarbeiter:} Ingrid Schumacher, Dipl.-Inf.\\
Klaus-Dieter Schumacher, Dipl.-Inf. 
\end{huge}

\end{center}

% Inhaltsverzeichnis

\setcounter{page}{0}

\newpage 
\pagenumbering{roman} 


\tableofcontents 

\newpage

% Start der eigentlichen Arbeit

\setcounter{page}{0}
\pagenumbering{arabic} 

\section{Ziel der Arbeit}
Die Android App \textbf{MobileTrainTestsoftware} dient als mobile Erweiterung der Leitstelle und ergänzt die auf den Steuerrechnern installierte Testsoftware.
Die momentan auf den Steuerrechnern installierte Testsoftware hat den Nachteil, dass zwei Personen zum Testen der jeweiligen Ebene notwendig sind.
Unmittelbar nach dem Schalten einer Komponente ist diese bereits geschaltet und es ist nicht leicht erkennbar, ob dies problemlos abgelaufen ist.
Mit der von uns entwickelten Android App MobileTrainTestsoftware kann eine geschaltete Komponente direkt beobachtet werden, da die Nutzung der App mobil möglich ist.

\subsection{Anforderungen}
Zu Beginn des Projektes wurde von unserer Auftraggeberin Frau Schumacher festgelegt, welche Anforderungen unsere App erfüllen soll.
Die zu entwickelnde App soll alle Ebenen der Modelleisenbahnanlage darstellen und steuern können.
Weichen und Streckenblöcke sollen mit dem Gleisplan dargestellt sein und durch interaktives Feedback steuerbar sein.
Außerdem soll die App über eine angenehme und moderne UI verfügen.

\section{Einleitung}\label{einleitung}
Hier folgt noch eine Einleitung zu der Modelleisenbahnanlage.


\section{Analyse bisheriger Technologien}
\subsection{Netzwerkkommunikation}
Die Kommunikation wird in der Modelleisenbahnanlage auf unterschiedliche Weise umgesetzt.
Dabei befinden sich alle Geräte im selben Netzwerk, über welches die Kommunikation läuft.

Zum einen gibt es die Leitstelle (Pi400). 
Sie kommuniziert mit auf der Anlage verteilten Steuerrechnern (RPi4) zum Ändern der Zustände von Weichen oder Streckenblöcken.
Die Leitstelle sendet dazu zunächst die eigene IP-Adresse per UDP-Broadcast an alle, im Netzwerk vorhandenen, Steuerrechner. 
Gleichzeitig wird ein Thread zum Annehmen einer TCP-Verbindung zu den jeweiligen Steuerrechnern gestartet. 
Jede TCP-Verbindung kann zum Senden von Informationen (neue Komponentenkonfiguration) oder dem Empfang von Zuständen der Komponenten verwendet werden. 
Dazu läuft jede TCP-Verbindung in einem eigenen Thread. Wenn eine definierte Anzahl Steuerrechner eine TCP-Verbindung zu der Leitstelle aufgebaut haben, wird ein Initialzustand (für Weichen und Streckenblöcke) an alle Steuerrechner gensendet. 
Eine Zustandsänderung einer Komponente wird per UDP übertragen und es wird auf eine Antwort des zuständigen Steuerrechners gewartet. 
Nachzulesen ist dies in der main.py und Network\_new.py in der Leitstellen Software.
Die Steuerrechner wiederum warten auf einen ankommenden UDP-Broadcast der Leitstelle. 
Nachdem die UDP- und TCP-Sockets eingerichtet wurden, wird der UDP-Socket abgehört. 
Erfasst der Steuerrechner ein UDP-Broadcast, so baut der Steuerrechner eine TCP-Verbindung zu der IP-Adresse der Leitstelle auf, welche in der UDP Nachricht mitgesendet wird.
Wichtig ist, dass der Steuerrechner eine TCP-Verbindung zum Port 6005 öffnet.
Würde an der Leitstelle ein anderer Port für die TCP-Verbindung geöffnet, so könnte keine Verbindung hergestellt werden.
Die IP-Adresse muss für den Steuerrechner nicht im String der UDP-Nachricht enthalten sein, da diese schon im Daten-Paket vorhanden ist.
Somit erwarten die Steuerrechner als Verbindungsaufbau UDP-Nachricht keinen spezifischen String.
Der Steuerrechner des Lift-Blocks erwartet allerdings einen spezifischen String.
Wird dieser nicht gesendet, so baut er keine Verbindung auf. 
Über die aufgebaute TCP-Verbindung wird eine Nachricht mit dem erfolgreichen Verbindungaufbau an die Leitstelle zurückgeschickt. 
Anschließend wird die UDP-Verbindung vom Steuerrechner in einem eigenen Thread abgehört. 
Nachdem alles initialisiert ist, wird in die Main-Loop übergegangen, wo das jeweils älteste Element einer zuvor initialisierten Queue behandelt wird. 
Wird eine Nachricht via UDP empfangen, so wird in dieser Nachricht überprüft, ob der jeweilige Steuerrechner angesprochen wird. 
Sollte das nicht der Fall sein, wird die Nachricht verworfen. 
Sollte das Level dem Level des Steuerrechners entsprechen, so wird der Befehl in der Queue abgelegt. 
Nach der Ausführung des Befehls, wird mit Hilfe des Communicators eine Bestätigung der Ausführung über TCP an die Leitstelle zurückgesendet. 
Nachzulesen ist dies in \_\_main\_\_.py und communicator.py in der Software des Steuerrechners.



\section{Umsetzung der Android App}

\subsection{Planung der Android App und Entwicklungsschritte}
Nachdem die Anforderungen an unsere App festgelegt wurden, haben wir eine Einführung in die Modelleisenbahnanlage bekommen. 
Dabei haben wir uns die Gleispläne und die eingesetzten Geräte (Steuerrechner, Leitstelle, Tablet) angesehen.
Zu der Steuerrung via Steuerrechner haben wir mit direkten I2C Befehlen einige Weichen testweise über einen der Steuerrechner geschaltet.
Dies half uns, ein Verständnis für die Befehle und Schaltungslogik der Anlage zu bekommen.
Als Resultat haben wir einen Arbeitsplan entwickelt.

Im Vorhinein haben wir vereinbart, Schritte arbeitsteilig anzugehen.
Yannic entschied sich, die Gestaltung der App sowie das Frontend zu übernehmen, während Thorben die Netzwerkkommunikation und das Backend übernahm.
Der Arbeitsplan enthielt folgende Schritte:
\begin{enumerate}
    \item Bezeichnung der Weichen und Streckenblöcke auf der Anlage mit Etikitiergerät anbringen
    \item Skizze für den Aufbau der App entwickeln (Yannic)
    \item Analyse der vorhandenen Technologie und Kommunikationstechnologie der App festlegen (Thorben)
    \item Mockup der App erstellen (Yannic)
    \item Erstes Element der Anlage in Praxisversuch schalten (Thorben)
    \item Alle anderen Funktionen umsetzen (Yannic und Thorben)
    \item Front- und Backend zusammenfügen
    \item Iteration ab Punkt 6, falls der Auftraggeber nicht zufrieden ist
    \item Abgabe des finalen Produkts und Dokumentation
\end{enumerate}

Nach der Analyse der vorhandenen Technologien haben wir uns bezüglich der Netzwerkkommunikation zuerst Gedanken darüber gemacht, wie wir Nachrichten an Steuerrechner senden können.
Unsere erste Überlegung war, alle Steuerrechner einzeln direkt via SSH anzusprechen.
Allerdings wurde uns schnell bewusst, dass diese Idee nicht umsetzbar ist, da die Steuerrechner eine dynamsiche IP-Adresse haben.
Für eine direkte Kommunikation via SSH bräuchten wir die IP Adresse der jeweiligen Steuerrechner.

Nachdem wir diese Idee aufgegeben hatten, haben wir uns an eine Internetprotokoll basierte Kommunikation wie zwischen Leitstelle und Steuerrechner festgelegt.
Um einen ähnlichen Ansatz wie bei der vorhandenen Kommunikation zu verfolgen, war es nötig, die bereits vorhanden Umsetzungen zu analysieren.
Nachdem wir dies getan haben, haben wir in einem ersten Test versucht festzustellen, ob die Kommunikation auch mit einer Android App funktioniert.
Getestet haben wir dies zu Beginn mit einem einfachen Button, der lediglich eine einzige Weiche schaltet konnte.
Die Steuerrechner sind beim Testen zu Beginn abgestürzt, da unsere Testversion noch keinen TCP-Listener hatte.
Nachdem die Steuerrechner allerdings korrekt mit der Leitstelle verbunden ware, war es kein Problem eine UDP-Nachricht von einem zweiten Gerät neben der Leitstelle an die Steuerrechner zu senden.
Das Ergebnis dieses Tests war, dass wir ohne Initialisierung der Steuerrechner eine Weiche korrekt schalten konnten.
Nach diesem Ergebnis haben wir die grafische Benutzeroberfläche entworfen und an einer TCP-Listener-Methode gearbeitet.

\subsection{Netzwerkkommunikation}
Die Android App befindet sich im selben WLAN wie die anderen Komponenten der Modelleisenbahnanlage.
Dies ist das Netzwerk \textbf{eisenbahn}, welches aktiv ist, sobald dies an der Fritz-Box eingeschaltet ist.
Die App hat im Android-Manifest die Erlaubnis, Netzwerkkommunikation zu nutzen.

Zuerst wird ein TCP-Listener gestartet, der den TCP-Port 6005 abhört.
Sollte der TCP-Listener einen Verbindungsaufbau erfassen, so wird dieser in einem neuen Thread behandelt.
Dies ermöglicht, eine TCP-Verbindung mit mehreren Clients herzustellen.
In diesem Fall sind die Clients die Steuerrechner der jeweiligen Modellbahn Ebenen.
Anders als bei der Leitstelle warten wir nicht auf eine Verbindung zu allen Steuerrechnern der Anlage, sondern öffnen so viele Threads, wie Verbindungen eingehen.

Über einen Button links unten kann man steuern, ob ein Verbindungaufbau zu den Steuerrechnern erfolgen soll.
Beim Öffnen der App ist unsere App im Disconnected Modus, was auch auf dem Button angezeigt wird.
Möchten wir eine Verbindung herstellen, so aktivieren wir den Button.
Es erscheint der Text connecting und der Button ist orange.
Im Hintergrund wird nun eine Nachricht via UDP-Broadcast an alle Netzwerkgeräte im selben Netzwerk gesendet.
Die Nachricht enthält einen String und das zugehörige Datenpaket, in dem auch die IP-Adresse enthalten ist.
Der UDP-Broadcast läuft wie auch der TCP-Listener in einem eigenen Thread.

Sobald ein Steuerrechner den UDP-Broadcast entdeckt, so stellt er eine TCP-Verbindung mit der Android App auf Port 6005 her.
Der Button, der vorher noch connecting angezeigt hat, ist nun grün und zeigt verbunden an.
Der Steuerrechner sendet nach dem erfolgreichen Herstellen der Verbindung einen String mit der Ebene, die er steuert (Ebene 0,1,2).
Die nun verbundenen Steuerrechner warten jetzt auf Nachrichten der Android App.

Wird eine Aktion wie etwa das Schalten einer Komponente in der App durchgeführt, so wird ein String mit dem dazugehörigen Befehl an die Steuerrechner via UDP-Nachricht gesendet.
Alle Steuerrechner empfangen demnach die Nachricht aber wie bereits bei der Analyse der bereits vorhandenen Technologien beschrieben, verarbeiten sie nur Nachrichten, die für sie vorgesehen sind.
Sollte ein Befehl einem Steuerrechner zugehörig sein, so setzt er den Befehl um und antwortet mit einem Acknologement, welches eine Kopie des von der App gesendete Strings ist.
Unsere App empfängt den String des Steuerrechners mit dem TCP-Listener und zeigt das Ack an.

Soll die Verbindung mit den Steuerrechnern unterbrochen werden, so kann über den Button mit einem Klick die TCP-Verbindung zu allen Clients geschlossen werden.
Nach dem Klick auf den Button, ändert sich das Aussehen von grün und der Anzeige verbunden zu rot und Disconnected.

\subsection{Architektur der App}

\section{Evaluation}


\end{document}
